\documentclass[a4paper,11pt]{article}
\usepackage[utf8]{inputenc}
\usepackage[T1]{fontenc}
\usepackage{geometry}
\geometry{margin=2.2cm}
\usepackage{xcolor}
\usepackage{graphicx}
\usepackage{tikz}
\usetikzlibrary{arrows.meta, positioning, shapes.geometric, fit, backgrounds}
\usepackage{everypage}
\usepackage{background}
\usepackage{amsmath}
\usepackage{booktabs}
\usepackage{float}
\usepackage{caption}
\usepackage{multirow}
\usepackage{enumitem}
\usepackage{array}
\usepackage{lmodern}
\usepackage[colorlinks=true, linktoc=all,
            linkcolor=sebgreen, urlcolor=sebgreen]{hyperref}
\captionsetup{font=small, labelfont=bf}
\setlength{\parskip}{0.7em}
\setlength{\parindent}{0pt}

% --- SEB colours ---
\definecolor{sebgreen}{RGB}{0,122,61}
\definecolor{sebdark}{RGB}{26,26,46}
\definecolor{lightgrey}{RGB}{244,244,244}
\definecolor{midgrey}{RGB}{180,180,180}
\definecolor{darkgrey}{RGB}{85,85,85}

% --- Page border ---
\backgroundsetup{
  scale=0.96, color=sebgreen, opacity=0.7, angle=0,
  position=current page.center,
  contents={
    
\begin{tikzpicture}[remember picture, overlay]
      \draw[line width=4pt, color=sebgreen]
        (current page.south west)
        rectangle (current page.north east);
    \end{tikzpicture}
  }
}

% --- SEB logo position (adjust these to move the logo) ---
\newcommand{\logoxshift}{-0.55cm}   % negative = move left from right edge
\newcommand{\logoyshift}{0.82cm}    % positive = move up from bottom edge
\newcommand{\logowidth}{1.8cm}      % logo size
\newcommand{\logoopacity}{0.6}      % 0.0 = invisible, 1.0 = fully solid

% --- SEB logo bottom-right on every page except cover ---
\newcounter{pagecounter}
\AddEverypageHook{
  \stepcounter{pagecounter}
  \ifnum\value{pagecounter}>1
    \begin{tikzpicture}[remember picture, overlay]
      \node[anchor=south east, xshift=\logoxshift, yshift=\logoyshift, opacity=\logoopacity]
        at (current page.south east)
        {\includegraphics[width=\logowidth]{logo.png}};
    \end{tikzpicture}
  \fi
}

% --- Header/footer ---
\usepackage{fancyhdr}
\pagestyle{fancy}
\fancyhf{}
\renewcommand{\headrulewidth}{0pt}
\fancyfoot[C]{\color{darkgrey}\small \thepage}

% --- Section colours ---
\usepackage{titlesec}
\titleformat{\section}{\large\bfseries\color{sebgreen}}{}{0em}{}[\color{midgrey}\rule{\linewidth}{0.6pt}]
\titleformat{\subsection}{\normalsize\bfseries\color{sebdark}}{}{0em}{}

% --- Table helpers ---
\newcolumntype{L}[1]{>{\raggedright\arraybackslash}p{#1}}

\begin{document}

% =============================================================================
% COVER
% =============================================================================
\begin{titlepage}
\thispagestyle{empty}

\begin{tikzpicture}[remember picture, overlay]
  \fill[sebgreen] (current page.north west) rectangle
    ([yshift=-4cm]current page.north east);
\end{tikzpicture}
\vspace*{3.5cm}
\begin{center}
{\Huge\bfseries\color{sebgreen} Probability of Default Model}\\[0.8em]
{\large\color{darkgrey}\itshape SEB Quantitative Risk Case}\\[2em]
{\normalsize\color{darkgrey} Magnus Foldeström}
\end{center}
\vspace{1.4cm}
\begin{center}
  \includegraphics[width=10cm]{logo.png}
\end{center}
\vspace{1.5cm}
\end{titlepage}

\clearpage
\tableofcontents
\clearpage

% =============================================================================
\section{Task and Data}
% =============================================================================

\textbf{Task:} Estimate the 12-month probability of default for each borrower.
Output is a continuous probability, not a binary decision.
Maximum 10 features. Compare at least two modelling approaches.

\textbf{Data:} 70{,}975 observations, 44 features, 2008--2013.
Default rate 1.81\%. One anomalous row (PERIOD = 1956) removed.
Temporal split: train 2008--2011 (53{,}468 rows, 884 defaults),
test 2012--2013 (17{,}506 rows, 403 defaults).

% =============================================================================
\section{Assumptions}
% =============================================================================

\begin{itemize}[leftmargin=1.5em]
  \item \textbf{Missing at random (MAR).} The 19 missing values across 18 rows
        are assumed to be MAR. Median imputation from training data is therefore valid.

  \item \textbf{Approved borrowers only.} The training data covers only borrowers
        who were approved under the previous policy. Rejected applicants were never
        observed, so we cannot know how they would have performed. If the new policy
        approves borrower profiles that were previously rejected, the model's
        predictions for those profiles are extrapolations beyond its training distribution.
\end{itemize}

% =============================================================================
\section{Methodology}
% =============================================================================

\textbf{Class imbalance.} The training set has a default rate of 1.65\%
(884 defaults out of 53{,}468 rows). Addressed at three levels:
(1) Weight-of-Evidence (WoE) transformation works with proportions within each class, inherently
robust to imbalance;
(2) Brier score is benchmarked against the naive baseline;
(3) \texttt{scale\_pos\_weight} for tree models is tuned via cross-validation.

\textbf{Preprocessing.} Only 18 of 70{,}975 rows had missing values (19 in total).
Dropping these rows would discard all other valid information they contain just to
avoid filling a single value. Median imputation was used rather than mean, because
distressed borrowers tend to produce extreme financial ratios that would skew a
mean-based fill.

\textbf{Feature selection pipeline:}

\begin{table}[H]
\centering
\begin{tabular}{L{3cm} r r L{8cm}}
\toprule
\textbf{Stage} & \textbf{In} & \textbf{Out} & \textbf{Criterion} \\
\midrule
Information Value (IV) filter        & 44 & 37 & IV $< 0.02$ dropped \\
Variance Inflation Factor (VIF) removal      & 37 & 29 & VIF $> 10$ dropped, iteratively \\
Forward stepwise & 29 & 10 & $50\%$ Gini $+ 50\%$ Brier (5-fold cross-validation (CV)) \\
\bottomrule
\end{tabular}
\caption{Feature selection funnel.}
\end{table}

\textbf{WoE transformation.} Each selected feature is replaced by its
Weight-of-Evidence value:
\[
  \text{WoE}_{\text{bin}} = \ln\!\left(\frac{\text{share of defaults in bin}}
                                             {\text{share of non-defaults in bin}}\right)
\]
Fitted on training data only.

\textbf{Models trained:}

\begin{table}[H]
\centering
\begin{tabular}{L{4.5cm} L{2.5cm} L{7.5cm}}
\toprule
\textbf{Model} & \textbf{Input} & \textbf{Purpose} \\
\midrule
Logistic Regression (LR/WoE) & WoE features & Primary model. Interpretable. Calibrated via Maximum Likelihood Estimation (MLE). \\
Logistic Regression (L1/WoE) & WoE features & Validation: confirms all 10 features carry independent signal. \\
XGBoost                      & Raw features & Benchmark. \texttt{scale\_pos\_weight=1} (CV-tuned). \\
LightGBM                     & Raw features & Benchmark. Leaf-wise boosting. Same tuning as XGBoost. \\
\bottomrule
\end{tabular}
\caption{Models trained and their purpose.}
\end{table}

\subsection{Method Justification}

\begin{table}[H]
\centering
\begin{tabular}{L{3.5cm} L{11cm}}
\toprule
\textbf{Method} & \textbf{Why it was chosen} \\
\midrule
Temporal split &
  A random split would leak future information since the same borrower can
  appear across multiple years. Out-of-time validation better reflects real deployment. \\[0.4em]
Median imputation &
  Robust to extreme values. Distressed borrowers tend to have extreme financial
  ratios that would pull a mean-based fill in the wrong direction. \\[0.4em]
IV screening &
  $\text{IV} = \sum_{\text{bins}} (\text{default share} - \text{non-default share})
  \times \text{WoE}$. Standard univariate screen in credit risk. Two practical
  reasons to include it: (1) banking convention, regulators expect it as part of
  a documented PD pipeline; (2) computational efficiency, IV is cheap to compute
  and reduces the candidate pool from 44 to 37 features before the more expensive
  cross-validated forward selection loop. \\[0.4em]
VIF removal &
  Detects multicollinearity involving combinations of features simultaneously,
  more comprehensive than pairwise correlation filters. \\[0.4em]
WoE transformation &
  Handles non-linear relationships, outliers (via edge bins), and class imbalance
  (works with proportions, not raw counts). Enables interpretable LR coefficients. \\[0.4em]
Forward stepwise &
  Jointly optimises discrimination (Gini) and calibration (Brier) in 5-fold CV.
  Stops when no remaining feature improves the combined criterion. \\[0.4em]
Logistic Regression &
  Interpretable coefficients (log-odds per feature). Calibrated via MLE. Faster
  and cheaper to train than tree models. Each borrower prediction can be fully
  decomposed into feature contributions ($\text{WoE} \times \text{coefficient}$). \\[0.4em]
XGBoost \& LightGBM &
  Trained as benchmarks to quantify the performance cost of choosing LR over a
  more flexible model. Without them, the claim that LR is ``good enough'' has no
  reference point. The gap is 7.1\% Gini when each model is evaluated on its own
  optimal feature set, which is a meaningful difference and an honest basis for
  the model choice discussion. \\
\bottomrule
\end{tabular}
\caption{Justification for each methodological choice.}
\end{table}

% =============================================================================
\section{Results}
% =============================================================================

All models evaluated on the out-of-time test set (2012--2013).
Gini $= 2 \times \text{AUC} - 1$ where AUC is the Area Under the Curve (discrimination).

LR and tree models ran separate forward stepwise selections so each model type
was evaluated on its own optimal feature set.

\begin{table}[H]
\centering
\begin{tabular}{L{6.5cm} @{\hspace{1cm}} L{6.5cm}}
\toprule
\textbf{LR features (WoE)} & \textbf{Tree features (raw)} \\
\midrule
\multicolumn{2}{l}{\textit{Shared by both models:}} \\[0.2em]
SUM\_ANM                      & SUM\_ANM \\
REMINDER\_6MTH\_CNT           & REMINDER\_6MTH\_CNT \\
CNG\_KUNDFORD                 & CNG\_KUNDFORD \\[0.4em]
\multicolumn{2}{l}{\textit{Unique to each model:}} \\[0.2em]
DELINQUENT\_TIMES\_18MTH\_CNT & SOLIDITET \\
CNG\_BEK                      & SATILLG \\
UTDELN                        & REVOLVING\_CREDIT\_PCT \\
SKULDRANTA                    & LIKV\_MEDEL\_OMS \\
AVK\_EGKAP                    & AVK\_TOTKAP \\
ARBKOST\_OMS                  & REMINDER\_24MTH\_CNT \\
BALRES                        & PERIOD\_WITH\_BANK \\
\bottomrule
\end{tabular}
\caption{Features selected by forward stepwise for LR (via LR internally) and
tree models (via XGBoost internally). Order does not indicate importance.
Three features were selected by both models.}
\end{table}

Brier score:
\[
  \text{BS} = \frac{1}{n}\sum_{i=1}^{n}(y_i - \hat{p}_i)^2
\]
Lower is better. Naive baseline:
\[
  \text{BS}_{\text{naive}} = p(1-p) = \frac{403}{17506}\times\frac{17103}{17506} = 0.0225
\]
where $p$ is the test-set default rate (2.30\%).

\begin{table}[H]
\centering
\begin{tabular}{lrrrrrr}
\toprule
\textbf{Model} & \textbf{Gini Train} & \textbf{Gini Test} & \textbf{Gini Gap} & \textbf{Brier Train} & \textbf{Brier Test} & \textbf{Brier Gap} \\
\midrule
LR (WoE)       & 0.9340 & 0.8898 & 0.044 & 0.0123 & 0.0178 & 0.0054 \\
LR (L1/WoE)    & 0.9340 & 0.8901 & 0.044 & 0.0123 & 0.0177 & 0.0054 \\
XGBoost        & 0.9986 & 0.9611 & 0.038 & 0.0029 & 0.0116 & 0.0087 \\
LightGBM       & 0.9989 & 0.9582 & 0.041 & 0.0028 & 0.0118 & 0.0090 \\
\midrule
Naive baseline & ---    & ---    & ---   & ---    & 0.0225 & ---    \\
\bottomrule
\end{tabular}
\caption{Model performance on the out-of-time test set.}
\end{table}

All models beat the naive Brier baseline. Tree models have a slightly smaller
train-test Gini gap (0.038--0.041) than LR (0.044), though the absolute test Gini
for trees is substantially higher (0.961 vs 0.890). Top-decile capture: if the 10\% riskiest borrowers are flagged for review, LR
captures 83\% of all actual defaulters within that group; XGBoost captures 93\%.
Of the 1{,}750 borrowers in that top decile, 19\% actually defaulted under LR
(334 out of 1{,}750), compared to 93\% capture meaning roughly 375 of the 403
total defaulters are found. The remaining defaulters are spread across the lower
9 deciles.

\begin{figure}[H]
\centering
\includegraphics[width=\textwidth]{cap_calibration.png}
\caption{Left: Cumulative Accuracy Profile (CAP). X-axis: fraction of borrowers sorted
by predicted risk (high to low). Y-axis: fraction of all defaults captured. A curve
closer to the top-left captures more defaults with fewer borrowers reviewed, corresponding
to a higher Gini.
Right: calibration curves. Each point shows the actual default rate in a bin against
the mean predicted probability for that bin:
$y = \text{defaults in bin} / \text{borrowers in bin}$.
A model on the diagonal is perfectly calibrated.}
\end{figure}

\subsection{PSI (Population Stability Index): Feature Stability}

\[
  \text{PSI} = \sum_{\text{bins}}
  \bigl(\text{test share} - \text{train share}\bigr)
  \times \ln\!\left(\frac{\text{test share}}{\text{train share}}\right)
\]

Thresholds: PSI $< 0.10$ stable; $0.10$--$0.25$ some shift; $> 0.25$ major shift.
All 10 LR features are stable between training and test periods. PSI was computed
for LR features only; tree models use a separate feature set not shown here.

\begin{table}[H]
\centering
\begin{tabular}{lrr}
\toprule
\textbf{Feature} & \textbf{PSI} & \textbf{Status} \\
\midrule
UTDELN                     & 0.0239 & Stable \\
SKULDRANTA                 & 0.0126 & Stable \\
ARBKOST\_OMS               & 0.0109 & Stable \\
SUM\_ANM                   & 0.0092 & Stable \\
CNG\_BEK                   & 0.0075 & Stable \\
BALRES                     & 0.0068 & Stable \\
REMINDER\_6MTH\_CNT        & 0.0058 & Stable \\
AVK\_EGKAP                 & 0.0048 & Stable \\
DELINQUENT\_TIMES\_18MTH\_CNT & 0.0036 & Stable \\
CNG\_KUNDFORD              & 0.0021 & Stable \\
\bottomrule
\end{tabular}
\caption{Population Stability Index for all selected features.}
\end{table}

% =============================================================================
\section{LR Coefficients}
% =============================================================================

WoE values range from negative (safer-than-average bins) to positive
(riskier-than-average bins). All coefficients are positive: predicted PD increases
as WoE increases. For any individual borrower, the contribution of each feature is
simply its WoE value multiplied by its coefficient.

\begin{table}[H]
\centering
\begin{tabular}{L{7cm} r r L{5.5cm}}
\toprule
\textbf{Feature} & \textbf{Coeff.} & \textbf{IV} & \textbf{Description} \\
\midrule
CNG\_KUNDFORD              & 1.729 & 0.094 & Change in accounts receivable (borderline IV) \\
BALRES                     & 0.991 & 0.577 & Retained earnings \\
SUM\_ANM                   & 0.886 & 2.877 & Sum of payment remarks (strongest by IV) \\
CNG\_BEK                   & 0.746 & 0.433 & Change in agreed overdraft limit \\
ARBKOST\_OMS               & 0.594 & 0.051 & Employee costs / turnover (borderline IV) \\
AVK\_EGKAP                 & 0.552 & 1.324 & Return on equity \\
REMINDER\_6MTH\_CNT        & 0.472 & 0.735 & Payment reminders last 6 months \\
DELINQUENT\_TIMES\_18MTH\_CNT & 0.394 & 0.836 & Times late on payments last 18 months \\
UTDELN                     & 0.307 & 1.931 & Dividend \\
SKULDRANTA                 & 0.263 & 0.430 & Average interest rate on debt \\
\bottomrule
\end{tabular}
\caption{LR coefficients and IV scores for all selected features.}
\end{table}

\textbf{CNG\_KUNDFORD} has the highest coefficient despite borderline IV.
IV is a univariate measure; low IV does not preclude a strong contribution in the
multivariate model. Forward selection uses Brier in cross-validation as part of its
criterion: if CNG\_KUNDFORD were merely predicting around the base rate, its WoE
values across bins would be near zero and Brier would not improve when adding it.
The Brier improvement in CV is the direct evidence that it carries genuine signal.

\textbf{L1 validation:} a logistic regression with L1 penalty ($C=0.1$) was trained
on the same features. No coefficient was driven to zero, confirming that all 10
features carry independent signal and that overfitting is not a concern.

\begin{figure}[H]
\centering
\includegraphics[width=\textwidth]{feature_importance.png}
\caption{Left: LR coefficients (WoE features). Right: XGBoost feature importance
by gain (tree features). The two models use different feature sets; features
appearing only in the LR panel have no XGBoost counterpart and vice versa.}
\end{figure}

% =============================================================================
\section{The Process}
% =============================================================================

A step-by-step account of how I approached this case from start to finish.

\begin{enumerate}[leftmargin=1.5em, label=\textbf{Step \arabic*.}]

  \item \textbf{Read the task brief carefully.}
        Noted the key requirements: probability output (not a binary decision),
        at most 10 features, at least two modelling approaches, and documented methodology.

  \item \textbf{Wrote down an overall approach before touching any code.}
        Planned the full pipeline upfront: temporal train/test split, IV screening,
        VIF removal, WoE transformation, forward stepwise selection, model training,
        and evaluation on both discrimination and calibration.

  \item \textbf{Read through varinfo.csv and reviewed case\_data.csv.}
        Spent time understanding what each variable actually measures before any
        modelling. This informed later decisions about which features made intuitive
        sense as credit risk predictors.

  \item \textbf{Considered whether a decision threshold should be determined.}
        Concluded the task only asks for a probability, not a lending decision.
        A cutoff requires bank-specific input on expected losses and risk appetite.
        Noted it as a natural next step but kept it outside scope.

  \item \textbf{Refreshed understanding of logistic regression.}
        Went back to the mathematical foundations to ensure full clarity on how
        coefficients, log-odds, and calibration interact before building the model.

  \item \textbf{Decided on a temporal train/test split.}
        Used 2008--2011 as training and 2012--2013 as test. A random split would
        leak future information since the same borrower can appear across multiple
        years. Out-of-time validation better reflects real deployment conditions.

  \item \textbf{Explored the data and removed one anomalous observation.}
        Found one row with PERIOD = 1956, clearly a data entry error. Removed it.
        Overall default rate: 1.81\% across all years, 2.30\% in the test period.

  \item \textbf{Handled missing values.}
        Only 18 of 70{,}975 rows had missing values (19 in total). Dropping them
        would discard all other valid information in those rows to avoid filling a
        single value. Chose median over mean because distressed borrowers tend to
        produce extreme financial ratios that would skew a mean-based fill.

  \item \textbf{Considered outlier treatment.}
        Decided against removing or clipping outliers. In credit risk, borrowers who
        default are often outliers by definition. WoE binning handles extremes
        automatically by placing them in the outermost bin.

  \item \textbf{Ran IV screening.}
        44 features reduced to 37 by dropping those with IV below 0.02. IV is the
        standard univariate screening metric in credit risk and directly linked to
        the WoE framework used downstream.

  \item \textbf{Removed the Pearson/Spearman correlation filter.}
        Had initially added a pairwise correlation filter between IV and VIF. After
        reflection, VIF already captures everything pairwise correlation does and
        more. Removing it improved LR test Gini from 0.868 to 0.890.

  \item \textbf{Ran VIF removal iteratively.}
        37 features reduced to 29 by dropping any feature with VIF above 10, one
        at a time, until all remaining features had VIF below the threshold.

  \item \textbf{Applied WoE transformation.}
        Fitted WoE bin edges on training data only and applied the same edges to the
        test set. This prevents any information from the test period influencing
        the transformation.

  \item \textbf{Ran forward stepwise selection.}
        29 features reduced to 10 using a combined criterion of 50\% Gini and 50\%
        Brier with 5-fold cross-validation. The combined criterion ensures the
        selected features contribute to both discrimination and calibration.

  \item \textbf{First full model run.}
        Trained all four models. XGBoost outperformed LR by roughly 4\% Gini.
        Considered the trade-off and concluded interpretability justifies recommending
        LR: each feature has a quantified, auditable contribution.

  \item \textbf{Tuned scale\_pos\_weight for tree models via cross-validation.}
        Ran a CV search over [1, 5, 10, 20, 40, 60] using the combined Gini+Brier
        criterion. Optimal value was spw=1 for both XGBoost and LightGBM, meaning
        no upweighting of the minority class improved the combined criterion.

  \item \textbf{Considered L1 and L2 regularisation.}
        L2 on LR was redundant given that forward selection already controls which
        features enter. L1 was retained as a validation tool: if it drove any
        coefficient to zero, it would suggest forward selection included a redundant
        feature. No coefficient was driven to zero, confirming all 10 features are necessary.

  \item \textbf{Computed PSI (Population Stability Index) for all selected features.}
        Measured how much each feature's distribution had shifted between the training
        period (2008--2011) and the test period (2012--2013). All 10 features fell
        below the 0.10 stability threshold. This step is primarily a documentation
        and regulatory requirement: forward stepwise selection already implicitly
        filters out unstable features, since a feature with inconsistent distribution
        across time would not produce consistent CV improvements and would not be
        selected. PSI confirms and formalises what the selection process already
        implied.

  \item \textbf{Reviewed the report for accuracy.}
        Went through the report and adjusted statements, removed anything that could
        not be backed up directly from the data, and made sure all methodological
        choices were correctly motivated.

  \item \textbf{Verified all task requirements.}
        Re-read the task brief systematically. Probability output: confirmed.
        At most 10 features: confirmed (exactly 10). At least two models: confirmed
        (four trained). Documented methodology: confirmed.

  \item \textbf{Wrote the report in \LaTeX.}
        Produced a final typeset report with properly rendered formulas, tables,
        and figures.

\end{enumerate}

% =============================================================================
\section{Further Work}
% =============================================================================

\begin{itemize}[leftmargin=1.5em]

  \item \textbf{Decision threshold.} The model outputs probabilities. Setting a
        cutoff for approve/decline requires expected loss inputs from the business.

  \item \textbf{Ongoing monitoring.} PSI and Gini should be tracked on new cohorts.
        Retrain if performance degrades materially.

  \item \textbf{Reject inference.} The training population covers approved borrowers
        only. The previous model that rejected applicants used different criteria,
        so we cannot know how our model would have scored them. This is a blind spot
        that cannot be resolved without additional data.

  \item \textbf{Individual explanations.} For LR with WoE, borrower-level explanations
        are already built in: the contribution of each feature is its WoE value
        multiplied by its coefficient. If XGBoost or LightGBM were chosen as the
        primary model instead, SHAP (SHapley Additive exPlanations) values would be needed.

  \item \textbf{Economic regime / concept drift.} The calibration curve shows that
        the model systematically underestimates default probabilities on the 2012--2013
        test set. A likely explanation is that the training and test periods represent
        two fundamentally different economic regimes. Training covers the acute phase
        of the 2008 financial crisis: sharp shock, elevated early defaults, then rapid
        stabilisation. The test period coincides with the European sovereign debt
        crisis: slow, prolonged stress with deteriorating macro conditions and defaults
        arising in companies that appeared stable under the previous regime. The model
        learnt to associate risk with acute, visible distress signals. When applied to
        a slow-burn debt crisis, it underestimates risk because the feature--default
        relationship has shifted and no macroeconomic variables are included to capture
        the regime change. This is classic concept drift. Recalibration on 2012 data
        or inclusion of macro features would partially address this. Platt scaling ---
        fitting a small logistic regression on top of the model's raw output to correct
        systematic over- or underprediction --- is a common ad hoc remedy, though it
        does not resolve the underlying regime shift.

  \item \textbf{Periodic retraining.} As new labelled cohorts become available, the
        model should be retrained to capture evolving credit behaviour.

\end{itemize}

% =============================================================================
\section{Anticipated Questions}
% =============================================================================

\textbf{Why choose LR over XGBoost despite lower Gini?}

XGBoost achieves 7.1\% higher Gini when each model is evaluated on its own optimal
feature set, but LR is preferred for two reasons.
First, credit decisions must be explainable at the borrower level: each feature's
contribution is its WoE value multiplied by its coefficient, giving a fully auditable
prediction. Second, LR is faster and cheaper to train and score,
which matters when the model is retrained periodically on new data.

\textbf{What does calibration mean in practice?}

A calibrated model means that among all borrowers assigned a predicted PD of 3\%,
approximately 3\% actually default. Banks use the predicted probability directly in
pricing (interest margins based on expected loss) and regulatory capital calculations.
If the model assigns 3\% when the true risk is 8\%, the bank underestimates its risk,
sets too low an interest rate, and holds insufficient capital. LR calibration via MLE provides this guarantee by construction. XGBoost has no
equivalent guarantee, though in this dataset it achieved a lower Brier score than
LR --- calibration should not be assumed to hold on future data without verification.

\textbf{The model was trained during the financial crisis. How does that affect usability?}

The model achieved strong discrimination on 2012--2013, but the calibration curve
reveals a systematic underestimation of default probabilities: the model is not
conservative. A likely explanation is that the training and test periods represent
two fundamentally different economic regimes. Training covers the acute phase of
the 2008 financial crisis: sharp shock, elevated early defaults, then rapid
stabilisation. The test period coincides with the European sovereign debt crisis:
slow, prolonged stress with deteriorating macro conditions and defaults in companies
that appeared stable under the previous regime. The model learnt to associate risk
with acute, visible distress signals. When applied to a slow-burn debt crisis, it
underestimates risk because the feature--default relationship has shifted and no
macroeconomic variables are included to capture the regime change. This is classic concept drift, and explains why Gini can remain high while
calibration deteriorates simultaneously. Retraining on more recent data is the proper remedy. Recalibrating via Platt scaling
--- fitting a small logistic regression on top of the model's raw output to correct
systematic over- or underprediction --- is a common ad hoc fix, but does not resolve
the underlying regime shift.

\textbf{Can the same borrower appear in both train and test, causing leakage?}

The test outcomes (whether a borrower defaulted in 2012--2013) are never visible
during training, so there is no leakage of future information. However, a related
problem exists within the CV loop used for forward stepwise selection. Because
CV folds were assigned randomly rather than temporally, a borrower's 2011 data
could end up in the validation fold while their 2010 data is used for training.
This is backwards in time: the model trains on the future and validates on the past.
More concretely, if a borrower defaulted in 2010, that signal could influence
training before the model is validated on that same borrower's 2009 data.
Temporal CV folds, where training always precedes validation in time, would
eliminate this. This is a known limitation of the current implementation.

\textbf{How do you know the model actually identifies defaults rather than just predicting the base rate?}

Two pieces of evidence. First, all models beat the naive Brier baseline of 0.0225,
which is the score of a model that always predicts the base rate for everyone.
A model predicting zero for all borrowers would be penalised heavily by the loss
function for each actual default. Second, LR achieves Gini 0.890, meaning it
meaningfully separates high-risk from low-risk borrowers. A model predicting
the same value for everyone would have Gini = 0.

\textbf{CNG\_KUNDFORD and ARBKOST\_OMS have borderline IV but were selected. Should you have raised the IV threshold?}

IV is a univariate measure. A feature with low IV can still contribute information
that no other feature in the model covers. Forward stepwise captures this through
the multivariate Brier and Gini improvement in cross-validation. If these features
were merely predicting around the base rate, their WoE values would be near zero
and CV Brier would not improve when adding them. The CV improvement is the direct
evidence of genuine signal. Forward stepwise is the principled filter; IV is a
fast pre-screen, not a strict selection criterion.

\textbf{How would you set a decision threshold in practice?}

A single threshold is likely too blunt. A more practical approach is to segment
borrowers into risk brackets based on predicted PD, for example five to ten rating
buckets, and set pricing, credit limits and approval rules per bucket. This gives
the bank flexibility to adjust its risk appetite per segment without changing the
model. The appropriate bucket boundaries depend on the bank's risk appetite and
the economics of each loan type, information that was not available in this case.

\textbf{If you could redo this from scratch, what would you change?}

Two things. First, avoid spending time on methods that were later removed because
they did not improve the model. Second, use temporal CV folds in forward stepwise
selection to eliminate within-training leakage.

\textbf{How do you know LR is not also overfitting, just less than XGBoost?}

The train-test Gini gap for LR is 0.044, which is small in absolute terms, and
the test Gini of 0.890 is strong for a 10-feature model. Additionally, the L1
validation showed that no coefficient was driven to zero, suggesting the model
is not carrying redundant features that inflate in-sample performance. Some degree
of overfitting is always possible, and further regularisation could be explored,
but the evidence available does not suggest it is a material concern.

\textbf{How often should the model be retrained?}

That requires empirical study and depends on factors such as the macroeconomic
environment and how quickly borrower behaviour evolves. Rather than following a
fixed schedule, PSI and Gini on new cohorts are the natural signals to monitor.
If PSI approaches 0.10 or Gini drops materially on a new cohort, that is a
concrete signal that retraining is warranted.

\textbf{Is 6 years of data enough for a stable PD model?}

The Basel convention is generally a minimum of 5 years of observation data.
The training set covers 2008--2011, which is 4 years. Whether the convention
applies to the training period alone or to the full dataset used for the final
model (2008--2013, 6 years) is something to verify against the specific
regulatory guidelines. If it applies to the final model, the requirement is met;
if it applies to the training period only, it represents a minor weakness.

\textbf{Why are the PSI values suspiciously low given economic turbulence?}

It is possible that the selected features measure structural properties of
companies that are less sensitive to economic cycles than one might expect.
It is also worth noting that this is a case study dataset, which may have been
constructed to be tractable. Whether the low PSI reflects genuine stability or
an artefact of the data is something that would require further investigation
in a real deployment setting.

\textbf{Why report Gini if Brier is sufficient for both model comparison and credit decisions?}

For an individual credit decision, the predicted probability is all that is needed.
For model comparison, Brier is in principle sufficient. However, Brier mixes ranking
and calibration into a single number, making it harder to diagnose where a model
fails. Two failure modes illustrate why both metrics are useful:

\begin{itemize}[leftmargin=1.5em]
  \item \textbf{Good Brier, low Gini.} PD estimates are close to the true rate on
        average but nearly identical for all borrowers. The model is well-calibrated
        but cannot separate high-risk from low-risk. Useless for credit decisions.
  \item \textbf{High Gini, poor Brier.} The model ranks borrowers correctly but the
        probability levels are wrong. Good structure, poor calibration. Useful for
        ranking but not for pricing or capital requirements.
\end{itemize}

Gini isolates ranking, Brier captures both ranking and calibration. Together they
give a more complete picture than either alone. Gini is also more stable across
portfolios with different default rates, since Brier is sensitive to the population
default rate. It is additionally the established metric that regulators and credit
analysts expect to see.

\textbf{Why 5 folds and not 10 or leave-one-out?}

Leave-one-out is computationally prohibitive with 53{,}468 training rows. Five-fold
is industry standard and a pragmatic balance: each fold contains roughly 10{,}000
rows, which is large enough for stable CV estimates. Ten folds would give slightly
lower bias in the CV estimate but higher variance and double the computation cost,
which is difficult to justify during development. The number of folds can be
increased before production if computation allows.

\textbf{How generalizable is the model to other markets?}

Probably not very. Credit behaviour, accounting standards and macroeconomic
conditions differ between countries. The WoE bins and coefficients are calibrated
to the Swedish credit landscape during 2008--2013, and there is no guarantee that
the same patterns hold elsewhere. A model intended for a different market would
need to be trained on local data.

\textbf{How does the model handle extreme outlier values?}

WoE binning places extreme values in the outermost bin together with other extreme
observations. A borrower with a debt-to-income ratio of 500\% and one with 80\%
may fall in the same bin and receive the same WoE value, so a single extreme
observation does not disproportionately influence the prediction. Tree models also
handle outliers naturally through their split-based structure.

\textbf{Have you considered whether any features contain future information?}

The features are assumed to be temporally correct since this is a case study dataset.
In production this would need to be verified carefully. For corporate loans, features
are derived from the most recent published annual report rather than self-reported
figures, since audited accounts are more reliable. This creates a timing risk: a loan
decision made in January 2011 should be evaluated using the 2009 annual report, the
latest available at that point, but if the database matches on fiscal year rather than
publication date, the 2010 figures may be incorrectly attached to that decision. The
model then trains on information that did not exist at the time of the decision, which
can make it appear more predictive than it would be in production. This is a
well-known problem in credit risk, and there is an entire area within machine learning
called point-in-time data management dedicated to ensuring that each data point only
contains information that was actually available at the specific decision date.

\textbf{Could WoE binning itself introduce overfitting?}

WoE binning is a discretisation that reduces the model's degrees of freedom. Instead
of fitting to every individual value of a feature, the model fits to a small number
of bins. This increases bias slightly but reduces variance, a deliberate bias-variance
tradeoff that favours generalisation. If the model were overfitting, we would expect
to see larger gaps between training and test Brier and Gini. The observed gaps are
small, which is consistent with the variance-reducing effect of binning.

\textbf{VIF only measures linear dependence. Does that cause problems for tree models?}

VIF is the correct tool for LR, which is a linear model where multicollinearity
directly destabilises coefficients. For XGBoost, which handles correlated features
differently through its tree-building process, a Spearman-based filter would have
been more appropriate. In practice, tree models ran their own forward stepwise
selection using XGBoost internally, so each model type was evaluated on its own
optimal feature set.

\textbf{What happens when Gini and Brier point in opposite directions during forward stepwise?}

The two metrics are normalised to the same scale and combined into a single criterion
with equal weight. A feature that improves Gini more than it worsens Brier produces a
positive net contribution and is added. This is precisely why the two metrics are
combined: to avoid selecting features that optimise one dimension at the expense of
the other.

\textbf{Is 884 defaults in the training set enough to learn a stable pattern?}

WoE transformation addresses class imbalance by working with proportions within each
class rather than raw counts, so the absolute number of defaults is less critical than
in raw feature models. With 10 selected features, there are roughly 88 defaults per
feature. A common rule of thumb in credit risk is at least 20--30 defaults per bin
for stable WoE estimates, which is satisfied here. The tree models additionally
support tuning of \texttt{scale\_pos\_weight} to handle imbalance, though
cross-validation determined that no upweighting was needed in this case.

\textbf{Should the model account for loan size?}

The model estimates the probability of default, not the expected loss. Loan size is
relevant for computing expected loss (EL = PD $\times$ LGD $\times$ EAD, where LGD is Loss Given Default and EAD is Exposure at Default), but that
is a separate step the bank performs after the model delivers its PD estimate. Loan
size was not available in the dataset. If it were, it would have been worth testing
as a feature, since larger loans may attract different borrower profiles. Working
with the data available is a practical constraint, not a modelling flaw.

\textbf{Why median imputation rather than a more sophisticated method such as KNN?}

Median is simple, fast and robust to outliers. More importantly, only 18 of 70{,}975
rows had missing values. With such a small fraction, the choice of imputation method
has negligible impact on model performance. The added complexity of K-Nearest Neighbour (KNN) imputation
would not be proportionate to the problem.

\textbf{Does the speed advantage of LR matter if training happens offline?}

For a single loan application, scoring speed is negligible regardless of model. The
advantage becomes relevant when batch-scoring large portfolios, for example to compute
regulatory capital requirements or monitor risk across hundreds of thousands of
borrowers simultaneously. In that context, LR is materially faster than XGBoost.
Training speed also matters if the bank automates retraining triggered by PSI
monitoring, since faster retraining means updated models are available sooner.

\textbf{How would you explain a rejected loan to the borrower?}

With LR and WoE, each feature's contribution to the predicted PD is its WoE value
multiplied by its coefficient. This means it is straightforward to identify which
features drove the high PD for a specific borrower, for example payment remarks or
high interest burden, and communicate that clearly. This also satisfies regulatory
requirements: under GDPR (General Data Protection Regulation) and EU credit regulation, borrowers have the right to a
meaningful explanation of automated decisions. LR with WoE fulfils that requirement
directly without additional tools. XGBoost would require SHAP values or similar
post-hoc methods to produce equivalent explanations.

% =============================================================================
\section{Conclusion}
% =============================================================================

\textbf{Recommended model: Logistic Regression with WoE transformation.}
Gini 0.890 and Brier 0.0178 on the out-of-time test set. Well-calibrated via MLE.
Fully interpretable coefficients.

Tree models (XGBoost, LightGBM) with \texttt{scale\_pos\_weight} = 1 (CV-tuned)
achieve higher Gini (XGBoost: 0.961, LightGBM: 0.958) and lower Brier (0.012). They are not recommended as
the primary model because:
\begin{enumerate}[leftmargin=1.5em]
  \item credit decisions must be explainable at the borrower level, and LR
        coefficients directly quantify each feature's log-odds contribution;
  \item LR is calibrated via MLE, providing a mathematical guarantee of calibration
        that XGBoost lacks by construction --- though in practice XGBoost achieved
        a lower Brier score on this dataset;
  \item LR is faster and cheaper to train and score;
  \item the performance gap versus trees is 7.1\% Gini, an acceptable cost given
        the interpretability and calibration guarantees LR provides.
\end{enumerate}

All 10 features stable (PSI $< 0.10$). Periodic monitoring recommended.

\subsection*{Known Limitations}

\begin{itemize}[leftmargin=1.5em]
  \item IV and VIF were computed on the full training set before the CV loop. In a
        strictly correct implementation, IV and VIF would be recomputed inside each
        fold so that feature selection never sees the validation rows. As implemented,
        the features that entered forward selection were determined partly by data that
        later served as validation data within CV folds. In practice this introduces
        only a mild bias given the large training set.
  \item No decision threshold determined: requires bank-specific loss assumptions.
  \item Trained on approved applications only: untested on borrower profiles outside
        the historical approval range.
  \item Static model: performance may degrade if market conditions shift materially.
\end{itemize}

\subsection*{Appendix: Pseudocode Summary}

\begin{verbatim}
data = load("case_data.csv");  data = data[data.PERIOD != 1956]
train = data[PERIOD <= 2011];  test = data[PERIOD >= 2012]
impute missing with training medians (18 rows affected)
IV filter (keep IV >= 0.02):  44 -> 37 features
VIF removal (VIF > 10, iterative):  37 -> 29 features
Forward stepwise (50% Gini + 50% Brier, 5-fold CV):  29 -> 10 features
Fit WoE bins on training data; apply to test (same edges)
Tune scale_pos_weight via 3-fold CV over [1,5,10,20,40,60] -> spw=1
Train: LR(C=1), LR(C=0.1,L1), XGBoost(spw=1), LightGBM(spw=1)
Evaluate: Gini, Brier, PSI, top-decile capture
Output: lr.predict_proba(X_test_woe)[:,1]  ->  pd_predictions.csv
\end{verbatim}

\end{document}
